\section{Límites del método, implicaciones de ingeniería y preguntas abiertas}

Hasta aquí, los tres hallazgos ya cuentan con casos: representaciones abstractas que se reutilizan entre dominios, varias rutas que calculan en paralelo y objetivos futuros que restringen por adelantado el token actual. La última sección ya no agrega más «diagramas de mecanismos vistosos»; en su lugar, examina qué pueden demostrar estos resultados y qué no, y cómo cambian el diseño y la evaluación de sistemas.

\lecturefigure{slide-62-papers-and-methods.jpg}{Más casos y materiales de método: el artículo interactivo y el artículo metodológico de Circuit Tracing}{01:08:27}

\subsection{Cinco limitaciones que no pueden omitirse}

Primero, el replacement model es una aproximación, y el reconstruction error puede contener mecanismos clave. Segundo, las feature label las interpretan personas, y pueden no ser únicas o no ser estables. Tercero, la attribution es una influencia local prompt-local, no un programa global fijo. Cuarto, la investigación muestra casos exitosos; el artículo interactivo discute explícitamente que el método solo da resultados útiles en una parte de los prompts. Quinto, y es el punto que más se enfatizó en clase: la attention no está explicada, y la selección, el enrutamiento y la competencia podrían estar ocurriendo justamente ahí.

\begin{warningbox}{No trates el attribution graph como el «código fuente del modelo»}
Se parece más a un boceto local de mecanismos generado por un instrumento experimental: contiene pistas causales reales, pero también errores de sustitución, poda, agrupamiento y decisiones de quien interpreta. La forma correcta de usarlo es formular predicciones, hacer intervenciones, buscar contraejemplos e informar qué partes no se pueden explicar.
\end{warningbox}

\subsection{Del descubrimiento de mecanismos al diseño de sistemas}

Estos casos no ofrecen directamente un botón para «eliminar las alucinaciones», pero sí cambian la forma de plantear el problema de ingeniería. En lugar de ajustar solo un umbral de rechazo, se pueden medir por separado la evidencia known-entity, la evidencia de la respuesta y la competencia IDK/refusal; en lugar de confiar en lo que el modelo dice de sí mismo, se puede verificar con un verifier externo, la trace de las herramientas y las intervenciones contrafactuales; en lugar de tratar un forward pass fijo como el único presupuesto de cómputo, se puede permitir que la reflection, un recurrent loop o el adaptive compute agreguen verificaciones cuando la confianza es baja.

En el nivel del sistema también hay que separar las señales de mecanismo de las restricciones del producto. Aunque una feature interna pueda predecir errores, también puede desplazarse de una versión del modelo a otra y no puede convertirse directamente en una API de largo plazo; aunque una intervención funcione en los prompts experimentales, puede dañar otras capacidades. Una ruta más robusta consiste en convertir los descubrimientos de mecanismos en conjuntos de pruebas reproducibles, métricas de monitoreo e hipótesis de falla, y luego validarlos mediante evaluación fuera de línea, red teaming, pruebas de regresión y un despliegue gradual. La interpretabilidad aporta resolución diagnóstica, pero no sustituye la ingeniería de seguridad completa.

\teachervoice{01:10:05--01:11:33, el docente llama a «cómo evitar las alucinaciones» el problema de mil millones de dólares y dice explícitamente que no conoce la respuesta. Propone direcciones como una mejor calibration, hacer que el reasoning model reflexione, aumentar el cómputo de autoverificación y el recurrent/adaptive compute, y a la vez reconoce que estas opciones podrían sacrificar capacidad o creatividad.}

\begin{importantbox}{Los cuatro puntos de ingeniería que más vale la pena llevarse}
\begin{enumerate}
  \item Tratar las explicaciones del modelo como salidas por verificar, no como una interfaz de hechos internos;
  \item Monitorear la incertidumbre, la seguridad y la deriva de objetivos durante todo el proceso de generación, no solo revisar el prompt;
  \item Validar las hipótesis de mecanismo con intervention contrafactual y la reproducción de traces;
  \item Asignar cómputo variable a los tokens difíciles, para que el modelo busque información, reflexione o se detenga cuando lo necesite.
\end{enumerate}
\end{importantbox}

\subsection{Un flujo de auditoría de circuit-tracing reutilizable}

\begin{lstlisting}[language=Python,caption={Flujo de auditoría que va de un problema de comportamiento a hipótesis de mecanismo falsables}]
behavior = define_prompt_and_counterfactuals()
features = fit_sparse_replacement_model(behavior.activations)
graph = trace_prompt_specific_attribution(features)
hypotheses = label_prune_and_group(graph)
for hypothesis in hypotheses:
    intervene_in_base_model(hypothesis)
    test_predicted_behavior_change()
report_reconstruction_error_failures_and_scope()
\end{lstlisting}

Este pseudocódigo pone a propósito «report failures» como el último paso obligatorio. Tener solo grafos exitosos sin tasa de fallas produce discovery bias; tener solo features legibles sin intervenciones produce un sesgo narrativo; tener solo intervenciones sobre un único prompt sin un conjunto de contrafactuales hace que un efecto local casual se confunda con un mecanismo robusto.

Una auditoría real también debe guardar la semilla aleatoria, la versión del diccionario, el umbral de poda, la intensidad de la intervención y el prompt completo, para que otra persona investigadora pueda reproducir el mismo grafo y juzgar si la conclusión proviene en realidad del mecanismo del modelo o de los hiperparámetros del análisis.

Cuando distintos hiperparámetros dan rutas distintas, informar esa inestabilidad es en sí mismo un resultado: indica que la evidencia actual solo sostiene conclusiones de mecanismo de grano más grueso y que no permite explicar con seguridad hasta el nivel de un solo nodo o una sola arista.

\subsection{Preguntas abiertas}

La sesión de preguntas y respuestas dejó varias líneas de investigación realmente difíciles: cómo explicar la selección y el enrutamiento de la attention; cómo determinar si una feature es «la misma» en distintos prompts y modelos; cómo cuantificar la cobertura del graph respecto del comportamiento del modelo original; cómo extender los mecanismos a nivel de token a contextos largos, llamadas a herramientas y el agent loop; cómo distinguir las rutas necesarias, las rutas suficientes y las rutas redundantes de respaldo; y cómo lograr que el modelo, sin perder creatividad, calibre mejor el conocimiento de sus propios límites.

\teachervoice{01:11:33--01:11:57, el docente advierte que la «hallucination» misma no siempre está claramente definida: en un texto largo, en qué token empieza exactamente el error, y si el error es de hechos o una desviación del contexto, a menudo no tienen un límite único. Las métricas de investigación deben definir primero su objeto; no se puede tomar una queja difusa sobre el producto como etiqueta de mecanismo.}

\begin{knowledgebox}{Un informe de investigación debe presentar tres tablas a la vez}
\textbf{Tabla de Coverage}: qué prompts/nodos se explicaron; \textbf{tabla causal}: qué intervenciones cambiaron el comportamiento según lo predicho; \textbf{tabla de failure}: qué grafos son inestables, tienen errores grandes o no se pueden nombrar. Solo cuando las tres están presentes puede quien lee juzgar el alcance real de aplicación del método.
\end{knowledgebox}

\subsection{Resumen de la sección}

Circuit tracing lleva la pregunta «qué podría ocurrir dentro del modelo» a un nivel en el que se puede graficar, intervenir y refutar, pero no es un lector completo. El resultado más valioso hasta ahora es que da lugar a problemas de ingeniería más finos y a interfaces experimentales: distinguir representación, ejecución y selección; registrar el error; validar la causalidad local; y conservar la incertidumbre en las conclusiones.
